\renewcommand{\Anno}{1996}

\renewcommand{\NumeroProblema}{3}

\renewcommand{\Testo}{ 
Determinare le soluzioni intere dell'equazione
$$x^3+2y^3=4z^3.$$}

\renewcommand{\Risposta}{L'unica soluzione \`{e} $(0,0,0)$.}

\renewcommand{\Aiutino}{
Intanto $x$ deve essere pari, quindi $x=2a$. 
Andando a sostituire e semplificando si ha che anche $y$ deve essere 
pari, quindi $y=2b$. Analogamente si arriva a $z=2c$ e si nota che 
$(a,b,c)$ sono soluzioni della stessa equazione risolta da $(x,y,z)$. 
Questo dice che ogni soluzione non nulla produce una soluzione 
pi\`{u} piccola, ma allora \ldots 
}

\renewcommand{\SoluzioneUfficiale}{\TitoloSoluzione
Iniziamo dimostrando il seguente risultato preliminare: se $(x,y,z)$ 
\`{e} una soluzione intera dell'equazione proposta, allora $x$, $y$ 
e $z$ sono necessariamente numeri pari e la terna $(x/2,y/2,z/2)$ 
\`{e} ancora una soluzione intera della stessa equazione. 

Infatti, dall'equazione segue immediatamente che $x$ \`{e} un numero
pari.  Posto $x=2a$, sostituendo e semplificando un fattore 2 si
ricava che
$$4a^{3}+y^{3}=2z^{3}.$$
Pertanto anche $y$ deve essere pari: posto $y=2b$, 
sostituendo e semplificando si ricava che
$$2a^{3}+4b^{3}=z^{3}.$$
Di conseguenza anche $z$ deve essere pari: posto $z=2c$, 
sostituendo e semplificando si ricava che
$$a^{3}+2b^{3}=4c^{3}.$$
Questo mostra che la terna $(a,b,c)=(x/2,y/2,z/2)$ \`{e} soluzione
intera della stessa equazione.

Utilizzando il risultato preliminare mostriamo ora che $(0,0,0)$ \`{e}
l'unica soluzione intera.  Intanto \`{e} evidente che tale terna \`{e}
effettivamente una soluzione.  Supponiamo ora per assurdo che esista
una soluzione $(x,y,z)$ in cui le 3 variabili non sono tutte nulle.
Per quanto appena detto, $x$, $y$ e $z$ sono tutti pari e la terna
$(x/2,y/2,z/2)$ \`{e} ancora una soluzione intera.  Iterando il
procedimento si trova che anche $(x/4,y/4,z/4)$ \`{e} ancora una
soluzione intera, e cos\`{\i} via.  Tuttavia, a forza di dividere per
2, uno dei 3 elementi prima o poi diventer\`{a} dispari e avremo
cos\`{\i} una soluzione in cui $x$, $y$ e $z$ non sono tutti pari, il
che contraddice quanto dimostrato precedentemente.

Ne segue che $(0,0,0)$ \`{e} l'unica soluzione intera dell'equazione 
data.
}

\renewcommand{\SoluzioniAlternative}{}

\renewcommand{\Osservazioni}{\TitoloOsservazione
I lettori pi\`{u} esperti avranno riconosciuto nell'ultima parte della 
soluzione un'applicazione del \emph{principio della discesa infinita}. In 
alternativa, la dimostrazione si poteva concludere con il \emph{principio 
del minimo intero}, ragionando come segue. 

Supponiamo che esistano delle soluzioni in cui le 3 variabili non sono
tutte nulle e consideriamo una (o la) soluzione per cui la somma
$|x|+|y|+|z|$ risulta minima (questa esiste appunto per il principio
del minimo intero).  Ma noi sappiamo che $(x,y,z)=(2a,2b,2c)$ e che
$(a,b,c)$ \`{e} ancora una soluzione con 
$$|a|+|b|+|c|=\frac{|x|+|y|+|z|}{2}<|x|+|y|+|z|.$$

Questo contraddice la minimalit\`{a} di $(x,y,z)$ e quindi esclude
l'esistenza di soluzioni oltre quella banale $(0,0,0)$.  }


